\section{Porównanie czasu sortowania}
\label{sec:CzasSortowania}

Pierwszym i podstawowym kryterium wykorzystanym do oceny wydajności badanych algorytmów był czas wykonania sortowania. 
Analizie poddano warianty sekwencyjne oraz równoległe algorytmów Quick Sort, Merge Sort, 
Bucket Sort i~Sample Sort dla zbiorów danych o rozmiarach 1~000, 10~000, 100~000 oraz 1~000~000 elementów.

W pierwszej kolejności przeanalizowano wpływ rozmiaru danych na czas wykonania poszczególnych algorytmów. 
Jako punkt odniesienia przyjęto zbiór danych \texttt{random\_int}, zawierający losowo wygenerowane liczby całkowite. 
Pozwala to na porównanie badanych algorytmów w jednakowych warunkach, 
bez uwzględniania na tym etapie wpływu stopnia uporządkowania danych lub występowania duplikatów.

\begin{figure}[htbp]
    \centering
    \includegraphics[width=0.90\textwidth]
    {images/chapter6/execution_time/execution_time_vs_data_size_random_int.png}
    \caption{Średni czas wykonania wariantów sekwencyjnych w zależności od rozmiaru zbioru \texttt{random\_int}.}
    \label{fig:time-data-size-sequential}
    
    {\small Źródło: opracowanie własne.}
\end{figure}

Dla wszystkich badanych algorytmów zwiększenie rozmiaru zbioru skutkowało wzrostem czasu wykonania, 
przy czym różnice dla różnych algorytmów szczególnie stawały się widoczne dla największego badanego zbioru.

Zestaw danych obejmujący 1~000~000 losowych liczb całkowitych uzyskał najkrótszy średni czas wykonania dla algorytmu Bucket Sort, 
osiągając około 2,04~s. Nieco dłuższego czasu wymagał Quick Sort z czasem około 2,25~s. 
Merge Sort wymagał około 3,41~s, natomiast Sample Sort osiągnął najdłuższy czas spośród porównywanych algorytmów, równy 4,95~s.
Wyniki wskazują, że różnice pomiędzy badanymi implementacjami zwiększają się wraz ze wzrostem rozmiaru przetwarzanego zbioru.

W celu bezpośredniego porównania wariantów sekwencyjnych i~równoległych w tabeli~\ref{tab:best-execution-time} 
zestawiono czasy uzyskane dla 1~000~000 losowych liczb całkowitych. 
Dla każdego algorytmu przedstawiono czas wariantu sekwencyjnego oraz najlepszy czas uzyskany spośród badanych konfiguracji równoległych.

\input{tables/chapter6/best_execution_time}

Zestawienie pokazuje, że dla każdego z badanych algorytmów możliwe było uzyskanie czasu krótszego niż w wariancie sekwencyjnym, 
niemniej najlepsze wyniki osiągęto przy różnej liczbie jednostek wykonawczych. 
Największe przyspieszenie w stosunku do wykonania sekwencyjnego odnotowano dla Sample Sort i~wyniosło ono około 2,23. 
Nie oznaczało to jednak najkrótszego bezwzględnego czasu sortowania. 
Najlepszy wynik czasowy uzyskał Bucket Sort, dla którego przy wykorzystaniu 16 jednostek wykonawczych wyniósł on około 1,09~s.

Uzyskane wyniki wskazują, że zastosowanie wariantu równoległego może prowadzić do skrócenia czasu sortowania,
jednak wielkość uzyskanej poprawy jest zależna od badanego algorytmu oraz zastosowanej konfiguracji równoległości.